\input{preamble.tex}
\title{Cyclic markings and the Artin applications of geometric twisting}
\author{Human Compression project}
\date{Private candidate, 11 October 2026}
\begin{document}
\maketitle
\begin{abstract}
A cyclic permutation of a family of named generators gives a single
parabolic formula for the resulting semidirect product. This formula proves
all intersection and flag assertions for the involutory and cyclic Artin
constructions in Piedade--Tranchida, arXiv:2505.24662v1, Section~5.
We give mutual presentation maps, including the alternating-word calculation,
and distinguish the retained Artin parabolic theorem from the elementary
extension argument. Source presentation corrections are stated explicitly.
\end{abstract}

\section{The precise input}
For a group generated by a finite named set $S$, write
$A(R)=\langle r:r\in R\rangle$ for $R\subseteq S$, including
$A(\varnothing)=1$. We call this a parabolic marking if
\begin{equation}\label{eq:ip}
 A(R)\cap A(Q)=A(R\cap Q)
\end{equation}
and the cosets of all $A(R)$ have the finite Helly property: every finite
pairwise-intersecting collection has a common element. The names denote
distinct nonidentity elements. The standard coset system has maximal
parabolics $A(S\setminus\{s\})$.

Here are the classical inputs retained for an arbitrary Artin--Tits
marking. Standard parabolic presentations embed, their intersections obey
(\ref{eq:ip}), and their cosets have this Helly property. These are the
standard-parabolic results invoked in the source's Theorems~2.11--2.12:
van der Lek's parabolic embedding and the intersection/product results
cited there from Godelle--Paris. In particular, each singleton Artin
parabolic is infinite cyclic. For a Coxeter marking the corresponding
facts are the classical Coxeter parabolic theorem, and singleton
parabolics have order two. This note does not reprove those external
results or identify their proof costs with zero.
Precisely, the source locators are Bourbaki, \emph{Groupes et alg\`ebres
de Lie}, Chapters 4--6 (2006 edition), Chapter 4, \S1, Theorem 2
and Exercise \S1.1(b); H.~van der Lek, \emph{The homotopy type of
complex hyperplane complements} (1983), Chapter II, Theorem 4.13;
and E.~Godelle--L.~Paris, \emph{$K(\pi,1)$ and word problems for
infinite type Artin--Tits groups, and applications to virtual braid
groups}, Math.~Z.~272 (2012), Proposition 4.5. We retain these
interfaces as the primary source does, rather than claiming fresh
proof verification of the entire cited works.

For clarity, Helly is precisely the flag-transitivity input. A flag of
standard cosets is a pairwise-intersecting family, so Helly gives a common
representative and extends it to a chamber. Conversely, replace each
coset of an intersection of maximal parabolics by the corresponding
family of maximal-parabolic cosets. Pairwise intersections are preserved;
intersecting cosets of the same maximal parabolic coincide. Flag
transitivity supplies the common representative. Formula (\ref{eq:ip})
identifies all generated parabolics with these intersections.

\begin{lemma}\label{lem:product}
In a parabolic marking, for $H=A(R)$, $K=A(Q)$ and $L=A(P)$,
\[
 HK\cap HL=H(K\cap L).
\]
The complex of cosets of maximal parabolics, with simplices the families
having common intersection, is a flag complex. Its coset geometry is
residually connected; its panel of type $s$ has $|\langle s\rangle|$
chambers.
\end{lemma}
\begin{proof}
If $z\in HK\cap HL$, the left cosets $H,zK,zL$ meet pairwise:
the first meets each of the others and the latter two contain $z$.
Helly gives $h=zk=zl\in H$, with $k\in K,l\in L$; hence
$k=l\in K\cap L$ and $z\in H(K\cap L)$. The reverse inclusion
is immediate. The complex assertion is exactly Helly. Every standard
residue has acting group $A(R)$ and maximal parabolics
$A(R\setminus\{s\})$, $s\in R$. If $|R|\geq2$ they generate
$A(R)$: each named generator belongs to a child missing a different
name. Its incidence graph is therefore connected (successive
multiplications by parabolic elements give incidence paths). Flag
transitivity identifies every residue with a translate of one of these.
The Borel is $1$ by (\ref{eq:ip}), and the rank-one residue is
$\langle s\rangle/1$, proving the panel assertion.
\end{proof}

\section{One cyclic marking theorem}
Let $S=Y\sqcup L$ with
$L=\{x_0,\ldots,x_{k-1}\}$ and $k\geq2$. Suppose $(A,S)$ is a
parabolic marking and $\sigma\in\Aut(A)$ fixes every $y\in Y$ and
sends $x_j$ to $x_{j+1}$, with subscripts modulo $k$.
Then $\sigma^k=\id$ since it fixes every generator. Set
\[
 G=A\rtimes C_k=\langle A,t\mid t^k=1,\ t^{-1}at=\sigma(a)
 \ (a\in A)\rangle,
 \qquad T=Y\sqcup\{x_0,t\}.
\]
This notation fixes the conjugation convention. The quotient onto $C_k$
proves that $t$ has exact order $k$. The set $T$ generates $G$, since
$x_j=t^{-j}x_0t^j$.

For $R\subseteq T$ define a subset $E(R)$ of $S$ by
\[
 E(R)=\begin{cases}
 R,&t\notin R,\\
 R\cap Y,&t\in R,\ x_0\notin R,\\
 (R\cap Y)\cup L,&t,x_0\in R.
 \end{cases}
\]
In the first line $R$ is read inside $Y\sqcup\{x_0\}\subseteq S$.
\begin{theorem}\label{thm:cyclic}
For every $R\subseteq T$ the subgroup generated by $R$ is
\begin{equation}\label{eq:all-subgroups}
 G(R)=\begin{cases}
 A(E(R)),&t\notin R,\\
 A(E(R))\rtimes C_k,&t\in R.
 \end{cases}
\end{equation}
The marking $(G,T)$ has the intersection and coset Helly properties.
Consequently its standard coset system is flag-transitive, residually
connected and firm, and its coset complex is flag. The product identity
of Lemma~\ref{lem:product} holds for every three generated parabolics.
Its panels have size $k$ at $t$ and $|\langle s\rangle|$ at every
$s\in Y\sqcup\{x_0\}$.
\end{theorem}
\begin{proof}
If $t\notin R$, the assertion is the embedding of $A$ in $G$.
If $t\in R$ but $x_0\notin R$, it fixes every named generator in
$R\cap Y$, giving the second case. If both belong to $R$, conjugating
$x_0$ produces exactly the whole orbit $L$, and this subgroup of $A$
is invariant under $t$. This proves (\ref{eq:all-subgroups}) without
an order calculation or finiteness assumption on $A$.

If both subsets contain $t$, the two $A$-parts intersect by
(\ref{eq:ip}); the orbit $L$ is present in their intersection precisely
when $x_0$ belongs to both subsets. Thus their intersection is the
subgroup for $R\cap Q$. If neither contains $t$, use (\ref{eq:ip})
inside $A$. In the mixed case an element of the intersection has trivial
$C_k$ coordinate. The smaller support contains at most $x_0$ from $L$,
so intersecting its support with the other support again gives exactly
$E(R\cap Q)$. This proves full generated-subgroup intersection.

It remains to prove Helly, which does not follow from subgroup
intersection alone. The maximal parabolics just computed are
\begin{equation}\label{eq:maximal}
 P_t=A(Y\cup\{x_0\}),\quad
 P_{x_0}=A(Y)\rtimes C_k,\quad
 P_y=A(S\setminus\{y\})\rtimes C_k\ (y\in Y).
\end{equation}
The $A$-supports in the latter groups are invariant under $C_k$.
Take pairwise-intersecting left cosets of these groups. Repeated types
coincide. If a $P_t$ coset occurs, translate its $C_k$ coordinate to
$1$. All other cosets project onto $C_k$ and have representatives
$(a,1)$. Their fibers at coordinate $1$ are cosets of the indicated
standard parabolics in $A$. Two cosets of the latter kind intersect
if and only if their $A$ fibers intersect, because they are Cartesian
products with $C_k$; incidence with the $P_t$ coset likewise takes
place in this fiber. Thus the $A$ fibers meet pairwise and base Helly
gives a common representative. If no $P_t$ coset occurs, the same
argument applies at coordinate $1$ without translation. This proves
Helly for maximal parabolics. Replace a coset of any generated
parabolic by its maximal-parabolic cosets, using the intersection
property just proved, to get Helly for all of them. Finally apply
Lemma~\ref{lem:product}. Every named generator is nontrivial and the
panel sizes are at least two, proving firmness.
\end{proof}
This is the rank-one cyclic case of the orbit construction used in
Piedade--Tranchida's papers on twisting. Its proof is included in full
so that no earlier application manuscript is an unpriced mathematical
premise. It does not assert that an arbitrary Shephard marking has
these properties.

\section{The involutory presentation and the doubled braid}
Write $\W_m(u,v)$ for the alternating word of \emph{length} $m$,
starting with $u$. It is not the ordinary power $(uv)^m$.
Let $A$ be the Artin group on $Y\sqcup\{x,x'\}$, with Artin
relations $\mathcal R_Y$ on $Y$, identical labels $m_y$ on
$(x,y)$ and $(x',y)$, and label $m$ on $(x,x')$.
All labels lie in $\{2,3,\ldots\}\cup\{\infty\}$; an infinite
label means no relation. Interchanging $x,x'$ is an automorphism of
the presented group, fixes $Y$, and has order two because the names
are distinct by the retained Artin parabolic theorem.

\begin{lemma}\label{lem:braid}
If $t^2=1$, then for every $m\geq1$,
\[
 \W_m(x,txt)=(xt)^m t^{m\bmod2},\qquad
 \W_m(txt,x)=(tx)^m t^{m\bmod2}.
\]
Consequently the braid of length $m$ on $x,txt$ is equivalent to the
braid of length $2m$ on $x,t$.
\end{lemma}
\begin{proof}
For $m=2r$, the two words are $(xtxt)^r=(xt)^{2r}$ and
$(txtx)^r=(tx)^{2r}$. For $m=2r+1$, append respectively $x$ and
$txt$ to those words, obtaining $(xt)^{2r+1}t$ and
$(tx)^{2r+1}t$. Cancellation of the common final $t$ in the odd case
proves the equivalence. The alternating words of length $2m$ are
exactly $(xt)^m$ and $(tx)^m$.
\end{proof}

\begin{theorem}\label{thm:fold}
The semidirect product which interchanges $x,x'$ has presentation
\begin{align*}
 \langle Y,x,t\mid {}&\mathcal R_Y,\ t^2=1,\ [t,y]=1\ (y\in Y),\\
 &\W_{m_y}(x,y)=\W_{m_y}(y,x)\ (m_y<\infty),\\
 &\W_{2m}(x,t)=\W_{2m}(t,x)\ (m<\infty)\rangle.
\end{align*}
This is a Shephard presentation, and its named marking satisfies every
conclusion of Theorem~\ref{thm:cyclic}.
\end{theorem}
\begin{proof}
Map the displayed generators to the semidirect product. They satisfy
every relator: the last by Lemma~\ref{lem:braid}; the others follow
from the defining action and the original Artin relations. They generate,
since $x'=txt$. Conversely, in the displayed presentation set
$x'=txt$. Conjugating a mixed braid by $t$, which fixes $y$, gives
the braid on $x',y$. Lemma~\ref{lem:braid} supplies the braid on
$x,x'$. The relations on $Y$ are unchanged. Thus the original Artin
presentation maps into this group, and conjugation by $t$ implements
its transposition of the two leaves. The semidirect-product universal
property gives the inverse map. Both composites fix all named
 generators, so these maps are inverse. This establishes the exact
 group and marking, not merely a quotient action. The cyclic theorem
with $k=2$ now applies. An edge absent from the drawing has label two,
so all commutations $[t,y]=1$ remain part of the presentation.
\end{proof}

For type $D_n$ ($n\geq4$), take the chain $a_1,a_2,\ldots,a_{n-1}$
with consecutive label three, attach $a_n$ by label three to $a_2$,
and let all other distinct pairs commute. Interchange $a_1,a_n$.
The new marking $t,a_1,a_2,\ldots,a_{n-1}$ has vertex orders
$2,\infty,\ldots,\infty$, consecutive labels $4,3,\ldots,3$ and
all nonconsecutive commutations. It is exactly $B_n(2,\infty)$.
This proves source Propositions~5.12--5.13; the arbitrary labels above
prove Theorems~5.14--5.15, including every stated intersection and
flag-complex assertion. The $t$ panels have size two and all remaining
panels are infinite.

There are two literal presentation qualifications. On source page~46
the graph attaches $a_n$ to $a_2$, whereas the printed relator list
uses $a_3$ and does not consistently bound the chain index. The proof
on page~47 uses the graph's $a_2$ convention, as do we. Also the
$B_n(p,q)$ display on page~41 omits the nonconsecutive commutations;
they are required by Definition~5.1 and by the displayed graph. They
are explicitly retained here. These are corrected graph presentations,
not verification of the malformed relator lists as printed.

\section{Cyclic amalgamation}
Let $A_0$ be an Artin group on $Y\sqcup\{x\}$ and
$C=A_0(Y)$. The parabolic embedding theorem identifies $C$ with the
Artin group on the induced graph on $Y$. Form the amalgam $M$ of
$k\geq2$ copies of $A_0$ along this common $C$. Its presentation has
$Y$ and $x_0,\ldots,x_{k-1}$, the relations $\mathcal R_Y$, and the
same braid between $x_j$ and each $y\in Y$. There are no relations
between distinct $x_j$. This follows directly from the universal
property of the amalgam; the embeddings of each factor are the
classical normal-form theorem. The resulting presentation is itself
Artin, with all distinct leaf-pair labels infinite.

\begin{theorem}\label{thm:amalgam}
Cyclic permutation of the copies defines an automorphism $\sigma$ of
$M$. With $t^{-1}x_jt=x_{j+1}$, its semidirect product has presentation
\[
 M\rtimes C_k\cong
 \langle Y,x_0,t\mid\mathcal R_Y,\ t^k=1,\ [t,y]=1\ (y\in Y),\
 \W_{m_y}(x_0,y)=\W_{m_y}(y,x_0)\ (m_y<\infty)\rangle.
\]
There is no relation on the pair $x_0,t$. This Shephard marking has
full generated-parabolic intersection and the product identity,
and its standard coset system is flag-transitive, residually connected
and firm with flag coset complex. Its panels have sizes $k$ and
infinity.
\end{theorem}
\begin{proof}
The permutation preserves the Artin relators and has order $k$ on
the distinct leaf names. The displayed generators map to the
semidirect product and generate it. In the converse direction define
$x_j=t^{-j}x_0t^j$. Since $t$ fixes $Y$, conjugation by $t^j$
transports every mixed braid to the required braid on $x_j,y$.
The relation $t^k=1$ closes the cycle. The Artin presentation of
$M$ and then its semidirect product therefore map into the displayed
group. Both composites fix $Y,x_0,t$, proving inverse maps. Apply
Theorem~\ref{thm:cyclic} to the Artin marking of $M$.
\end{proof}
These are source Theorems~5.16--5.17. The source's literal statement
writes $A_0\rtimes C_k$ where its construction and proof require
$M\rtimes C_k$. There need not even be the indicated cyclic action
on the single copy $A_0$. Its page~51 preliminary braid display also
has identical words on both sides; the right side must start with
$y$. When $k=2$, the order-two vertex is represented by no loop under
Definition~5.1, rather than by a loop labelled two. None of these
conventions affects the corrected presentation just proved.

\section{Free products and quotient actions}
If $A,B$ are Artin or Coxeter groups with disjoint standard markings,
$A*B$ has the union presentation, equivalently all cross labels are
infinite. Its generated subgroup on $R$ is
$A(R\cap S_A)*B(R\cap S_B)$ by reduced normal forms. Two such
subgroups intersect by intersecting the two factor markings: the
unique alternating reduced word belongs to both exactly when each
syllable does. The coset Helly theorem for free products proved in
\emph{Free products and converse restrictions in coset geometries}
then applies to the two classical input geometries. Its hypotheses
allow arbitrary Borel subgroups; here both Borels are trivial.
The resulting standard marking therefore satisfies
Lemma~\ref{lem:product}. This recovers source Theorem~5.4 and all
three clauses of Corollary~5.5, for arbitrary choices of the two
Artin/Coxeter factors. Each source panel size is preserved.

Proposition~\ref{prop:shephard} below proves the faithful graph action
asserted in source Lemma~5.6, including the distinctness step missing
from its short proof. Quotient constructions still retain their
specified indexed action even when subgroup images coincide.

For a group $A$ with a specified action of $B$, a $B$-invariant
subgroup $H$ has invariant normal closure $N$: every image of
$a^{-1}ha$ is a conjugate of an element of $H$, and applying the
inverse automorphism gives equality. For arbitrary $H$, the subgroup
$\langle b(h):b\in B,h\in H\rangle$ is invariant, by reindexing
$b$ after any further action. These prove source Lemmas~5.8--5.9.

The induced action on $A/N$ sends each image $A_iN/N$ to the image
indexed by $bi$. Keep this indexed action as part of the data even
if some subgroup images coincide. Every orbit-intersection condition
involving the acting group and named types is consequently unchanged.
This proves the admissibility assertion of Proposition~5.10 with
$\widetilde A_i=A_iN/N$, the well-defined image subgroup; the source's
notation $A_i/N$ need not be defined. No flag transitivity, intersection
property, firmness or residual connectedness is inferred for this
quotient. Those are separate hypotheses when applying twisting.

\section{Separating the named Shephard generators}\label{sec:reflection}
We now close the injectivity point in source Lemma~5.6 without assuming
an embedding theorem for general Shephard groups. This argument constructs
a representation separating the named generators; it does not assert that
the representation is faithful on the entire group.

\begin{proposition}\label{prop:shephard}
In the presentation of Definition~5.1, every named generator has its
prescribed order and distinct names give distinct elements. Consequently
the action of the labelled graph automorphism group on the presented
Shephard group is faithful.
\end{proposition}
\begin{proof}
Let $V$ have basis $(e_i)_{i\in S}$ over $\mathbb C$. For vertex order
$p_i<\infty$ choose $z_i=\exp(2\pi\mathrm i/p_i)$, and for infinite
order choose $z_i=2$. Thus $z_i\ne1$ in every case, and equal vertex
orders give equal $z_i$. Fix square roots $u_i^2=z_i$.
Define linear maps
\[
 R_i(e_i)=z_i e_i,\qquad R_i(e_j)=e_j+c_{ij}e_i\quad(j\ne i).
\]
Choose the two coefficients independently for each unordered pair,
after ordering its endpoints as $i<j$. For a commuting pair take
$c_{ij}=c_{ji}=0$; for label infinity do the same, since no braid
needs to be imposed. For finite labels greater than two take
\[
 c_{ij}=1,\qquad c_{ji}=
 \begin{cases}
 -4z_i\cos^2(\pi/m),&m\text{ odd},\\
 2u_i u_j\cos(\pi/k)-z_i-z_j,&m=2k.
 \end{cases}
\]
Proper labelling ensures $z_i=z_j$ in the odd case.
The maps are invertible, with determinant $z_i$. The geometric-series
formula for their exceptional row shows $R_i^{p_i}=1$ for finite $p_i$;
its eigenvalue $z_i$ gives exact order $p_i$. If $p_i=\infty$, the
eigenvalue two gives infinite order.

We verify the braid relations on the whole vector space. On
$V_{ij}=\langle e_i,e_j\rangle$, put
\[
 A=\begin{pmatrix}z&c\\0&1\end{pmatrix},\qquad
 B=\begin{pmatrix}1&0\\d&w\end{pmatrix},\qquad P=AB.
\]
For $m=2k$, the trace and determinant of $P$ are
$2uv\cos(\pi/k)$ and $(uv)^2$, where $z=u^2,w=v^2$.
Its distinct eigenvalues are $uv\exp(\pm\pi\mathrm i/k)$.
Thus $P^k=-(uv)^kI=(BA)^k$, the required braid.

For $m=2r+1$, put $z=w$, $q=-\cos(2\pi/m)$ and use
$cd=-4z\cos^2(\pi/m)$. Then $\operatorname{tr}P=2zq$ and
$\det P=z^2$. Let $U_{-1}=0$, $U_0=1$ and
$U_{j+1}(q)=2qU_j(q)-U_{j-1}(q)$. Cayley--Hamilton gives
\[
 P^r=z^{r-1}U_{r-1}(q)P-z^rU_{r-2}(q)I.
\]
Direct multiplication also gives
$PA-BP=(z+cd)(A-B)$. Since $1+cd/z=2q-1$, subtraction yields
\[
 P^rA-BP^r=z^r\bigl(U_r(q)-U_{r-1}(q)\bigr)(A-B)=0.
\]
For the last equality write $q=\cos\theta$ with
$\theta=\pi-2\pi/m$. The recurrence has solution
$U_j(q)=\sin((j+1)\theta)/\sin\theta$; hence the difference is
zero because $\cos(m\theta/2)=0$. The two sides before subtraction
are precisely the alternating words of length $2r+1$.

To pass from $V_{ij}$ to $V$, set
$D=(z-1)(w-1)-cd$. When $D\ne0$, the common fixed space of
$R_i,R_j$ is complementary to $V_{ij}$: solve its two fixed-vector
linear equations by adjusting the $e_i,e_j$ coordinates. Both maps
are the identity on that complement, so the verified pair relation
holds on $V$. The exceptional values $D=0$ cause no gap. In the odd
case vary $z$ while keeping all other row coefficients fixed and
$cd=-4z\cos^2(\pi/m)$; $D$ is a nonzero polynomial in $z$.
Every entry of the braid difference is a polynomial vanishing at all
but finitely many complex $z$, and thus vanishes also at the exceptional
values. In the even case vary $u$ with the nonzero $v$ fixed;
$D=u^2v^2+1-2uv\cos(\pi/k)$ is again a nonzero polynomial.
The same polynomial argument applies, including the finitely many
excluded zero values. For a commuting pair the off-diagonal pair
coefficients are zero and the two exceptional rows commute directly.
All defining relations therefore hold on $V$.

The presentation now gives a homomorphism to $\operatorname{GL}(V)$.
If $i\ne j$, the $i$th diagonal entry of $R_i$ is $z_i\ne1$,
whereas that of $R_j$ is one. Thus these matrices, and consequently
the named group elements, are distinct. A labelled graph permutation
and its inverse preserve all defining relators, so they induce inverse
group automorphisms. A nonidentity graph permutation moves a named
generator to a different one, proving faithfulness.
\end{proof}
The constructed representation is used only to separate names and
certify their orders. No parabolic embedding, intersection or coset
Helly property for arbitrary Shephard groups follows from it.

\section{Scope and attribution}
The Artin/Coxeter embedding and coset theorems, the free-product
normal-form theorem, and the already reviewed free-product Helly
argument are the precise retained inputs. The cyclic parabolic proof
and the mutual presentation maps are written here. The groups and
applications are Piedade--Tranchida's results; no priority claim is
made for this reorganization. Finite checks, if supplied, validate
word conventions and examples only and cannot replace the uniform
arguments. This candidate has not yet received distinct review and
claims neither complete source recovery nor full formalization.
\end{document}
